\documentclass[12pt, a4paper]{article}

% ── Paquetes ──────────────────────────────────────────────
\usepackage[utf8]{inputenc}
\usepackage[T1]{fontenc}
\usepackage[spanish]{babel}
\usepackage{geometry}
  \geometry{margin=2cm}
\usepackage{graphicx}
\usepackage{float}
\usepackage{hyperref}
  \hypersetup{colorlinks=true, linkcolor=black, citecolor=blue, urlcolor=blue}
\usepackage{listings}
\usepackage{xcolor}
\usepackage{amsmath}
\usepackage{booktabs}
\usepackage{caption}
\usepackage{parskip}
\usepackage{amssymb}
\usepackage{enumitem}
\usepackage{array}
\usepackage{tabularx}
\usepackage{pdflscape}
\newcolumntype{Y}{>{\raggedright\arraybackslash}X}
\usepackage{tikz}
\usetikzlibrary{positioning, shapes.multipart, arrows.meta}

\lstset{
  basicstyle=\ttfamily\small,
  keywordstyle=\color{blue}\bfseries,
  commentstyle=\color{gray},
  stringstyle=\color{orange},
  numbers=left,
  numberstyle=\tiny\color{gray},
  breaklines=true,
  frame=single,
  captionpos=b
}

% ══════════════════════════════════════════════════════════
\begin{document}

% ── CARÁTULA ──────────────────────────────────────────────
\begin{titlepage}
  \centering

  \includegraphics[width=0.2\linewidth]{img/logo-fai.png}
  \hfill
  \includegraphics[width=0.2\linewidth]{img/logo-unco.png}

  \vspace{0.5cm}
  {\large Departamento Teoría de la Computación}\\[0.2cm]
  {\large \textbf{Facultad de Informática}}\\[0.2cm]
  {\large Universidad Nacional del Comahue}\\[2cm]

  {\Huge \bfseries Laboratorio de inteligencia artificial}\\[0.5cm]
  {\Large Trabajo Práctico Obligatorio N\textsuperscript{o} 1}\\[0.3cm]
  {\LARGE \bfseries Parte 1: entendimiento del problema y trabajos relacionados}\\[0.3cm]
  {\large Predicción de la duración de desvíos de ruta en flotas de transporte de personal}\\[1cm]

  \begin{tabular}{lll}
    \toprule
    \textbf{Apellido y Nombre} & \textbf{Legajo} & \textbf{Email} \\
    \midrule
    Reibold, Sebastian Alejandro & FAI-3854 & sebastian.reibold@est.fi.uncoma.edu.ar \\
    Fernandez Gramajo, Bautista & FAI-3874 & bautista.fernandez@est.fi.uncoma.edu.ar \\
    Kissner, Davor & FAI-3185 & davor.kissner@est.fi.uncoma.edu.ar \\
    \bottomrule
  \end{tabular}

  \vfill
  {\large\quad Cursado 2026 \quad --- \quad 7 de noviembre de 2026}
\end{titlepage}

% ── ÍNDICE ────────────────────────────────────────────────
\tableofcontents
\newpage

% ══════════════════════════════════════════════════════════
\section{Introducción}

Este trabajo se enmarca en el convenio entre la empresa Tecfield SRL y la Facultad de Informática de
la UNComa, en la materia Laboratorio de Inteligencia Artificial. Tecfield cuenta con un detector
automático que compara la posición GPS de cada vehículo contra su hoja de ruta y registra un evento
cuando la distancia supera los 500 metros. En el instante en que aparece el primer punto fuera de
ruta, la sala de monitoreo no puede anticipar si el desvío durará minutos o más de una hora, lo que
dificulta priorizar la atención entre un volumen alto de alertas.

El proyecto busca estimar, en ese primer instante, la duración esperada del desvío. El trabajo
completo se presenta en tres documentos. Este primero plantea el problema, los objetivos y los
trabajos relacionados. El segundo describe los datos y la construcción del corpus. El tercero
presenta los modelos, las métricas y los resultados.

% ══════════════════════════════════════════════════════════
\section{Motivación}

Tecfield SRL opera flotas de vehículos dedicadas, principalmente, al transporte
de personal hacia obras y yacimientos de la región. Cada vehículo tiene asignada
una \emph{hoja de ruta} con un recorrido planificado, y un detector automático
compara periódicamente su posición GPS contra la geometría de ese recorrido.
Cuando la distancia supera un umbral (actualmente 500 metros), el detector
registra un \emph{evento de desvío} y agrupa las lecturas consecutivas que
permanecen fuera de ruta bajo un mismo identificador. El evento se cierra cuando
el vehículo retoma la ruta o cuando se alcanza el tope de cuatro horas del
detector, en cuyo caso no hay un retorno efectivamente medido.

El problema surge en el instante en que aparece el primer punto fuera de ruta
(que en adelante llamamos $T_0$). En ese momento la sala de monitoreo no tiene
forma de distinguir una separación momentánea ---el vehículo toma una calle
alternativa y retoma la ruta en uno o dos minutos--- de un evento prolongado que
puede extenderse durante una hora o más. Como el detector está en producción
desde 2026 y genera del orden de 100.000 lecturas por mes, el volumen de
alertas es alto y todas se atienden con la misma prioridad. Esto satura la
atención de los operadores y hace que los recursos de monitoreo no se destinen
necesariamente a los casos que más lo requieren.

Un agravante es que no existen etiquetas humanas confiables sobre el motivo de
cada desvío (falla de GPS, hoja de ruta mal configurada, vehículo mal asignado,
recorrido operativo distinto del plan), por lo que no es posible entrenar un
clasificador de desvíos ``relevantes'' versus ``irrelevantes''. Lo que sí queda
registrado de manera sistemática es \emph{cuánto duró} cada evento. Esto
convierte a la duración en una variable objetivo disponible, objetiva y
directamente útil para la operación.

Contar con una estimación temprana de la duración esperada de un desvío,
calculada en $T_0$ y sólo con información conocida hasta ese instante,
permitiría priorizar la atención de las alertas más críticas desde el momento
en que se generan, en lugar de esperar a que el evento se resuelva por sí mismo
o alcance el tope del detector. En otros términos, se trata de construir un
mecanismo de \emph{alerta temprana} (\emph{early warning}) sobre un evento que
ya fue detectado, pero cuya gravedad todavía se desconoce. La importancia de la
solución es que para Tecfield, reduce la carga de trabajo en monitoreo y mejora el tiempo de reacción ante desvíos largos.

% ══════════════════════════════════════════════════════════
\section{Objetivos}

\subsection{Objetivo general}

Desarrollar un prototipo que permita estimar, en el momento en que se detecta
el primer punto fuera de ruta de un vehículo ($T_0$), la duración esperada del
evento de desvío, de modo de apoyar la priorización de alertas en la sala de
monitoreo de Tecfield.

\subsection{Objetivos específicos}

\begin{enumerate}
  \item \textbf{Construir} un panel evento-a-evento a partir de las colecciones
  provistas por Tecfield (\texttt{desviosHojaRuta}, \texttt{hojaDeRuta} y
  \texttt{rastreoSatelital}), con la duración observada de cada evento como
  variable objetivo.
  \item \textbf{Analizar} el conjunto de datos construido para caracterizar la
  distribución de duraciones, la calidad y cobertura de las variables
  disponibles en $T_0$, y las limitaciones del corpus.
  \item \textbf{Desarrollar} un modelo predictivo que estime la duración del
  evento utilizando exclusivamente la información conocida hasta $T_0$.
  \item \textbf{Evaluar} el desempeño del modelo frente a referencias simples
  (\emph{baselines}) mediante una separación de los conjuntos de entrenamiento
  y prueba por vehículo, analizando el error según distintos segmentos de
  interés operativo.
\end{enumerate}
% ══════════════════════════════════════════════════════════
\section{Trabajos relacionados}

El problema se ubica en la intersección de tres líneas de
investigación bien establecidas: la predicción de desvíos respecto de rutas
planificadas, la detección de trayectorias anómalas a partir de GPS y la
predicción de la duración de incidentes viales. A continuación se describen
cuatro trabajos representativos y se analiza qué aporta cada uno al proyecto.

\subsection{Predicción de desvíos en rutas de última milla
  \cite{konovalenko2026}}

Se estudia sobre datos reales de una empresa de
logística noruega, si es posible predecir que un conductor se desviará de la
ruta planificada por el sistema de optimización y cuantificar la magnitud de
ese desvío. Disponen de la ruta planificada y de la ruta realmente ejecutada
para 19.647 rutas, 400 conductores y más de 10.000 ubicaciones, y el dataset
fue publicado. La hipótesis central es que los desvíos no son aleatorios sino
manifestaciones del \emph{conocimiento tácito} del conductor.

Se definen tres métricas de desvío entre ruta planificada y real (calidad de la
secuencia de paradas, distancia de edición normalizada y variación relativa de
la distancia total) y las usan para construir dos tareas: una regresión que
estima el grado de desvío y una clasificación binaria (desviada / no desviada)
según umbrales sobre las tres métricas. Las variables de entrada son
propiedades estáticas (identificador del conductor, día de la semana, país,
cantidad de paradas, experiencia reciente) y secuenciales (secuencia de
paradas, tipo de parada, si es depósito). Proponen una arquitectura de red
neuronal con \emph{embeddings} y un bloque secuencial intercambiable (LSTM,
atención o CNN), y la comparan contra cuatro \emph{benchmarks}: promedio
histórico global, promedio histórico por conductor y día, Random Forest y
perceptrón multicapa. La evaluación usa una partición temporal
(entrenamiento con rutas antiguas, prueba con las más recientes) estratificada
por conductor y día, con MSE, RMSE y MAE para regresión. Reportan mejoras del
9--19\,\% en regresión respecto de los \emph{benchmarks}.

\textbf{Relación con el proyecto.} Es el trabajo más cercano en cuanto a
contexto: ruta planificada versus ruta real, datos operativos de una empresa
y predicción de desvíos. Son directamente reutilizables (i) la metodología de
evaluación contra \emph{benchmarks} simples basados en
promedios históricos, (ii) el protocolo de agregar variables de manera
incremental y conservar sólo las que mejoran, y (iii) el conjunto de métricas
de regresión. La diferencia principal es el momento y la unidad de predicción:
ellos predicen \emph{antes} de que comience el viaje, a nivel de ruta completa,
si habrá desvío y cuánto se alterará la secuencia de paradas; nosotros
predecimos \emph{durante} el viaje, a nivel de evento ya detectado, cuánto
tiempo durará. Además, la variable objetivo de ellos es una similitud
acotada en $[0,1]$, mientras que la nuestra es una duración en segundos con
distribución muy asimétrica. La arquitectura secuencial con LSTM no aplica
directamente porque en $T_0$ no hay secuencia de paradas, aunque podría
considerarse si se incorpora la telemetría previa como serie temporal. También
difiere el criterio de partición: ellos separan por fecha y estratifican por
conductor, de modo que cada conductor aparece en los tres subconjuntos,
mientras que en nuestro caso se separa por vehículo justamente para impedir que
el modelo memorice el comportamiento de un móvil visto en entrenamiento.

\subsection{Predicción de la duración de incidentes viales
  \cite{grigorev2022}}

Se aborda la predicción de cuánto durará un
incidente de tránsito (accidente, avería, obstrucción) a partir de la
información disponible en el reporte inicial, sobre tres conjuntos de datos
de arterias y autopistas de Sídney y San Francisco. Proponen un marco
\emph{bi-nivel}: primero un clasificador decide si el incidente será corto o
largo, buscando el umbral que equilibra las clases y el desempeño; luego un
modelo de regresión específico para cada grupo estima la duración en minutos.
Incorporan un procedimiento de remoción de valores atípicos y una técnica de
optimización conjunta (IEO-ML) que extiende varios modelos base de aprendizaje
automático. Una conclusión relevante es que los incidentes cortos y largos
deben modelarse por separado, y que el enfoque propuesto supera a los modelos
base en la mayoría de los escenarios analizados.

\textbf{Relación con el proyecto.} Este trabajo resuelve un problema
estructuralmente idéntico al nuestro aunque en otro dominio: predecir la
duración de un evento en el momento en que se lo detecta, con información
parcial, sobre datos operativos con fuerte asimetría y valores extremos. Son
aplicables (i) la idea de tratar por separado los eventos cortos y largos,
posiblemente como una etapa de clasificación previa a la regresión, (ii) el
tratamiento explícito de \emph{outliers} en la variable objetivo, que se
corresponde con nuestra tarea de definir umbrales para eventos anómalamente
cortos o largos, y (iii) el uso de modelos basados en árboles y ensambles como
línea principal. La diferencia está en las variables de entrada: los reportes
de incidentes incluyen tipo de incidente, cantidad de carriles afectados,
condiciones climáticas y flujo de tráfico, mientras que nosotros contamos
principalmente con variables geoespaciales y de telemetría del propio
vehículo.

\subsection{Detección online de trayectorias anómalas
  \cite{chen2013}}

iBOAT es un método para detectar en tiempo real si
la trayectoria GPS de un taxi se aparta de las rutas ``normales'' históricas
entre un origen y un destino, e identificar qué tramo es responsable de la
anomalía. La ciudad se discretiza en una grilla, cada trayectoria se
representa como secuencia de celdas y, a medida que llegan nuevos puntos, se
calcula el soporte que tiene el tramo reciente en las trayectorias históricas:
un soporte bajo indica anomalía. Lo evalúan sobre las trayectorias de 7.600
taxis durante un mes, con desempeño cercano al máximo (AUC $\ge 0{,}99$), y
analizan unas 43.800 trayectorias anómalas, encontrando que la mayoría
corresponde a fraudes deliberados de los conductores.

\textbf{Relación con el proyecto.} iBOAT resuelve la etapa \emph{anterior} a
la nuestra: decidir si un vehículo está fuera de ruta. Tecfield ya tiene
resuelto ese paso con un detector propio basado en distancia a la geometría de
la hoja de ruta, por lo que no necesitamos reimplementarlo. El trabajo es útil
por dos motivos. Primero, muestra cómo representar y comparar trayectorias GPS
contra un patrón de referencia, lo que sugiere variables para $T_0$ (por
ejemplo, cuán ``frecuente'' es la zona donde se produce el desvío en el
historial del vehículo o de la obra). Segundo, ilustra que las anomalías
detectadas mezclan causas muy distintas (fraude, cortes de calle, rutas
alternativas legítimas), lo que refuerza la limitación de nuestro corpus:
los eventos de desvío incluyen tanto desvíos reales como falsos desvíos por
fallas de GPS o hojas mal configuradas, sin etiqueta que los distinga. La
diferencia central es que iBOAT es un problema no supervisado de detección,
mientras que el nuestro es supervisado de regresión.

\subsection{Modelos de duración basados en riesgo
  \cite{hojati2013}}

Se modela la duración de incidentes en
autopistas mediante modelos de supervivencia (\emph{hazard-based}), en los que
la variable de interés es el tiempo hasta que el evento finaliza y la función
de riesgo describe la probabilidad de que termine en el instante siguiente,
dado que aún no terminó. Ajustan distribuciones paramétricas (Weibull,
log-logística, log-normal) por tipo de incidente y estudian qué factores
aceleran o retrasan la finalización.

\textbf{Relación con el proyecto.} Este enfoque es relevante por una
característica particular de nuestros datos: los eventos que alcanzan el tope
de cuatro horas del detector no tienen un final observado, es decir, son
observaciones \emph{censuradas por la derecha}. Los modelos de supervivencia
están diseñados exactamente para incorporar ese tipo de casos en lugar de
descartarlos, lo que podría aprovechar información que un modelo de regresión
convencional pierde. En la primera etapa del proyecto se seguirá la regla
indicada en la consigna, que separa esos eventos por no corresponder a un
retorno medido; este trabajo queda como
referencia para una eventual extensión en la que se los trate como datos
censurados. También aporta la idea de que la distribución de duraciones suele
ser log-normal o Weibull, lo que justifica evaluar una transformación
logarítmica de la variable objetivo.

\subsection{Comparación}

La Tabla~\ref{tab:comparacion} resume los trabajos analizados frente a la
propuesta de este proyecto.

\begin{landscape}
\begin{table}[p]
  \centering
  \footnotesize
  \renewcommand{\arraystretch}{1.25}
  \begin{tabularx}{\linewidth}{>{\raggedright\arraybackslash}p{2.3cm} Y Y Y Y Y}
    \toprule
    & \textbf{Konovalenko 2026} & \textbf{Grigorev 2022} & \textbf{Chen 2013} & \textbf{Hojati 2013} & \textbf{Este proyecto} \\
    \midrule
    \textbf{Problema} & Predecir si una ruta será desviada y cuánto & Predecir duración de incidentes viales & Detectar en línea trayectorias anómalas & Modelar duración de incidentes & Predecir duración de un desvío ya detectado \\
    \textbf{Tipo de tarea} & Regresión + clasificación & Clasificación (corto/largo) + regresión & Detección no supervisada & Análisis de supervivencia & Regresión \\
    \textbf{Momento de predicción} & Antes de iniciar la ruta & Al reportarse el incidente & En cada nuevo punto GPS & Al reportarse el incidente & En el primer ping fuera de ruta ($T_0$) \\
    \textbf{Unidad} & Ruta completa & Incidente & Trayectoria origen--destino & Incidente & Evento (\texttt{grupoDesvioId}) \\
    \textbf{Datos} & 19.647 rutas planificadas y reales, 400 conductores & Reportes de incidentes de 3 redes viales (AU, EE.UU.) & GPS de 7.600 taxis, un mes & Reportes de incidentes en autopistas & Lecturas del detector (jul--ago 2026), hojas de ruta, GPS \\
    \textbf{Variables} & Conductor, día, país, largo, experiencia, secuencia de paradas & Tipo de incidente, carriles, clima, tráfico & Secuencia de celdas de grilla & Tipo de incidente, ubicación, hora & Distancia a ruta, velocidad, posición, hora, tipo de vehículo y hoja, telemetría previa \\
    \textbf{Modelos} & Redes con embeddings + LSTM/atención/CNN; RF y MLP como benchmark & Ensambles de árboles y otros modelos base; marco bi-nivel & Soporte de subtrayectorias sobre histórico & Weibull, log-logística, log-normal & Árbol de decisión y Random Forest \\
    \textbf{Evaluación} & Partición temporal estratificada; MSE, RMSE, MAE & Múltiples escenarios de validación; error en minutos & AUC contra trayectorias etiquetadas & Bondad de ajuste, análisis de factores & Partición por vehículo; MAE, RMSE, error mediano \\
    \textbf{Qué se toma} & Benchmarks históricos, selección incremental de variables, métricas de regresión & Separar cortos y largos; manejo de outliers; ensambles & Representación de trayectorias; conciencia de causas mixtas & Tratamiento de censura (tope de 4\,h); transformación log & --- \\
    \textbf{Qué no aplica} & Predicción previa al viaje; secuencia de paradas; su partición es temporal y la nuestra por vehículo & Variables de reporte no disponibles & Etapa de detección ya resuelta & Foco explicativo, no predictivo & --- \\
    \bottomrule
  \end{tabularx}
  \caption{Comparación entre los trabajos relacionados y la propuesta del proyecto.}
  \label{tab:comparacion}
\end{table}
\end{landscape}

En síntesis, ninguno de los trabajos resuelve el mismo problema, pero en
conjunto cubren los tres componentes que necesitamos: \cite{konovalenko2026}
provee la metodología experimental para desvíos de ruta con datos de empresa,
\cite{grigorev2022} y \cite{hojati2013} muestran cómo se aborda la predicción
de duración de eventos con distribuciones asimétricas y datos censurados, y
\cite{chen2013} enmarca el problema desde la detección sobre GPS y advierte
sobre la heterogeneidad de causas detrás de un desvío.

% ══════════════════════════════════════════════════════════
\begin{thebibliography}{9}

\bibitem{konovalenko2026}
A.~Konovalenko, L.~M. Hvattum y K.~A.~H. Iversen.
\newblock Predicting last-mile delivery route deviations using machine learning.
\newblock \emph{Expert Systems With Applications}, 298:129921, 2026.
\newblock \url{https://doi.org/10.1016/j.eswa.2025.129921}

\bibitem{grigorev2022}
A.~Grigorev, A.-S. Mihaita, S.~Lee y F.~Chen.
\newblock Incident duration prediction using a bi-level machine learning
  framework with outlier removal and intra--extra joint optimisation.
\newblock \emph{Transportation Research Part C: Emerging Technologies},
  141:103721, 2022.
\newblock \url{https://doi.org/10.1016/j.trc.2022.103721}

\bibitem{chen2013}
C.~Chen, D.~Zhang, P.~S. Castro, N.~Li, L.~Sun, S.~Li y Z.~Wang.
\newblock iBOAT: Isolation-based online anomalous trajectory detection.
\newblock \emph{IEEE Transactions on Intelligent Transportation Systems},
  14(2):806--818, 2013.

\bibitem{hojati2013}
A.~T. Hojati, L.~Ferreira, S.~Washington y P.~Charles.
\newblock Hazard based models for freeway traffic incident duration.
\newblock \emph{Accident Analysis \& Prevention}, 52:171--181, 2013.

\end{thebibliography}
\end{document}